\documentclass[11pt]{article}
\usepackage[a4paper,margin=1in]{geometry}
\usepackage{amsmath,amssymb,mathtools,bm}
\usepackage{enumitem,booktabs,array}
\usepackage{tcolorbox}
\usepackage{hyperref}
\usepackage[protrusion=true,expansion=false]{microtype}
\hypersetup{colorlinks=true,linkcolor=blue,urlcolor=blue,citecolor=blue}
\setlist[itemize]{topsep=2pt,itemsep=2pt}
\setlist[enumerate]{topsep=2pt,itemsep=2pt}
\newtcolorbox{keybox}{colback=blue!3!white,colframe=blue!40!black,boxrule=0.4pt,arc=1mm}
\newtcolorbox{alertbox}{colback=red!3!white,colframe=red!40!black,boxrule=0.4pt,arc=1mm}
\newtcolorbox{answerbox}{colback=green!3!white,colframe=green!40!black,boxrule=0.4pt,arc=1mm}
\newtcolorbox{drillbox}{colback=yellow!5!white,colframe=orange!60!black,boxrule=0.4pt,arc=1mm}
\newcommand{\EV}{\mathbb{E}}
\newcommand{\Var}{\mathrm{Var}}
\newcommand{\blank}[1]{\underline{\hspace{#1}}}

\title{W04-01: Gorton \& Metrick (2012, JFE) \\
\large ``Securitized Banking and the Run on Repo'' --- Active Recall Workbook}
\author{Banking \& Risk Management --- Exam Prep}
\date{\today}
\begin{document}\maketitle

\begin{keybox}
\textbf{Paper in one sentence.} The 2007--08 financial crisis was at heart a \emph{run on the repo market}: repo haircuts jumped from $\approx 0\%$ on non-subprime collateral in early 2007 to $\approx 45\%$ by late 2008, draining hundreds of billions of dollars of short-term funding from securitised banking, just as depositor runs did in 1907 or 1933.

\textbf{Placement.} The workhorse ``modern bank run'' paper. Reframes the GFC as a shadow-banking run rather than a housing-losses-on-solvent-banks story. Foundational for the Keys-et-al (W04-02), Mian--Sufi (W04-03) and Ivashina--Scharfstein (W04-04) papers.
\end{keybox}

\section*{Round 1: Complete Walk-Through}

\subsection*{1.1 Stylised facts of securitised banking}
\begin{itemize}
\item Prime brokers and dealers finance long positions in securitised products (ABS, MBS) via short-term \emph{repo}: they post collateral and borrow cash overnight or for short maturities.
\item Loan-to-value of repo is controlled by the \emph{haircut} $h$: cash advanced equals $(1-h)\times$ market value of collateral.
\item In calm times, haircuts on non-subprime collateral are essentially zero $\Rightarrow$ close to 100\% leverage on these positions.
\end{itemize}

\subsection*{1.2 The shock}
\begin{itemize}
\item ABX indices on subprime MBS start falling early 2007.
\item LIB--OIS spread widens; uncertainty about counterparty collateral quality rises.
\item Cash lenders (MMMFs, securities lenders) demand higher haircuts on \emph{all} securitised collateral, not just subprime.
\item The weighted-average haircut for the GM index rises to $\approx 45\%$ by late 2008.
\end{itemize}

\subsection*{1.3 Mechanics of a ``run''}
Each 1 pp increase in haircut reduces the amount of cash that any given collateral value can borrow. If repo books cannot shrink fast enough, dealers fire-sale collateral. This maps directly onto Diamond--Dybvig: depositors (cash lenders) run on demandable short-term liabilities (repo).

\subsection*{1.4 The GM repo-haircut index}
GM construct a weighted-average haircut across asset classes from dealer data. The index is the paper's signature contribution and has been used in many subsequent papers.

\subsection*{1.5 Empirical analysis}
GM regress a measure of ``risk'' (LIB--OIS, ABX index) against repo-haircut changes and show the two co-move tightly. They estimate that the run drained hundreds of billions in short-term funding.

\subsection*{1.6 Caveats}
\begin{alertbox}
\begin{itemize}
\item Data on repo haircuts are dealer-specific and GM stitch together from industry contacts; not public.
\item ``Run'' is stylised: it is information-sensitive withdrawal, not a classic bank run on demand deposits.
\item Debate over whether repo or tri-party repo or ABCP was the most important channel (see Krishnamurthy et al.~2014).
\end{itemize}
\end{alertbox}

\section*{Round 2: Level 1 Fill-in}
\begin{enumerate}
\item Repo mechanics: cash lender advances \blank{1cm} $\times$ collateral value.
\item GM haircut index rose from $\sim$\blank{1cm}\% to $\sim$\blank{1cm}\% between early 2007 and late 2008.
\item Repo run is a run on demandable \blank{2cm} liabilities.
\item Key stress indicators: \blank{2cm} spread and \blank{1cm} index.
\item GM's analogy: 2007--08 repo run $\leftrightarrow$ \blank{2cm} bank runs.
\end{enumerate}

\section*{Round 3: Level 2 Prompts}
\begin{enumerate}
\item Derive the leverage multiplier from a haircut: if haircut rises from $h_0$ to $h_1$, by how much must the repo book shrink?
\item Map the repo run to Diamond--Dybvig: who is the depositor, who is the bank?
\item Why did haircuts on \emph{non-subprime} collateral rise during the crisis?
\item Discuss the implications for LOLR policy: should the Fed backstop repo?
\end{enumerate}

\section*{Round 4: Minimal Scaffolding}
From a blank page: (i) repo mechanics and haircut, (ii) GM index construction, (iii) haircut spike numbers, (iv) Diamond--Dybvig analogy, (v) two caveats.

\section*{Round 5: Blank Page}
Essay: was the GFC really a run on repo? Evaluate GM against competing narratives.

\vspace{6pt}
\begin{drillbox}
\textbf{Speed Drill.} (1) Cash-to-collateral formula. (2) GM index peak. (3) Two stress indicators. (4) Who are the ``depositors'' in securitised banking? (5) One caveat.
\end{drillbox}

\section*{Quick Reference \& Answer Key}
\begin{answerbox}
\begin{itemize}
\item \textbf{Repo.} Cash $=(1-h)\times$ collateral value. Haircut $h$ is the margin.
\item \textbf{GM index.} $\sim 0\%$ early 2007 $\to \sim 45\%$ late 2008 on non-subprime collateral.
\item \textbf{Analogy.} Cash lenders $\leftrightarrow$ depositors; dealers $\leftrightarrow$ banks; repo $\leftrightarrow$ demand deposits.
\item \textbf{Information sensitivity.} Collateral becomes info-sensitive under doubt $\Rightarrow$ haircut spikes.
\end{itemize}
\end{answerbox}

\begin{keybox}
\textbf{30-second exam pitch.} Gorton--Metrick (2012) reframe the 2007--08 crisis as a run on securitised banking. Dealers and prime brokers finance long securities positions via repo; cash lenders control leverage with a haircut $h$. In calm times $h\approx 0$ on non-subprime collateral; during the crisis, the GM weighted-average haircut index rose to $\sim$45\%, draining hundreds of billions from dealer funding. This maps onto Diamond--Dybvig: repo is demandable, cash lenders are depositors, dealers are banks. The shock started with subprime ABX deterioration in 2007 but spread to all securitised collateral as information sensitivity rose. The paper is the canonical evidence that the GFC was a shadow-banking run, with clear policy implications for LOLR extensions to the shadow-banking sector.
\end{keybox}

\end{document}
